\chapter{Mountain Infrastructure, Validation and Risks / Infraestrutura de Montanha, Validação e Riscos}

\section*{English}
The station requires resilient power, UPS, batteries, generator, optical links, thermal management, lightning protection, EMC control, perimeter security and logistics. Site selection must consider terrain, horizon, local RFI, climate, access, energy and fibre availability. Environmental qualification must cover wind, vibration, rain, humidity, salinity, UV and thermal gradients.

The timing infrastructure includes at least two external GNSS antennas with surveyed coordinates and phase centres. They are mounted above or near the radome with clear sky visibility, physical separation from the VHF/UHF Yagis and filtered, shielded feeds to the timing room. Redundant multi-constellation receivers support cross-checking, fault detection and atomic-clock disciplining.

The civil interface is a reinforced-concrete parallelepiped base measuring approximately $4\,\mathrm{m}\times4\,\mathrm{m}$ in plan and $3\,\mathrm{m}$ in clear structural height. Its upper slab follows the lower cut plane of the radome at $35^\circ$ S, while the compact footprint and increased elevation preserve lower-angle antenna visibility without pointing sensors at the immediate ground. A service opening is provided in the front wall so personnel and equipment can enter the interior without removing RF panels.

The thermal solution is zoned. Heat generated by LNAs, ADCs, ASICs, FPGAs, optical transceivers and DC/DC converters is conducted to local cold plates and rejected through monitored air or liquid paths. The central timing and fusion core has an independent thermal zone. Condensation sensors, inlet/outlet temperature sensors and power telemetry support closed-loop control. Thermal phase drift is measured, timestamped and included in the calibration model.

The structural solution uses an internal load-bearing frame with sealed joints, node plates and gussets. Each external Yagi support/boom axis continues from the internal pyramid apex through the triangular face, while dielectric panels provide environmental protection and controlled RF transmission. Acceptance therefore covers static strength and dynamic stability: wind, vibration, modal response, thermal expansion, rain loading, maintenance access and repeated thermal cycles.

The development path is P0 system simulation, P1 three-node VHF/UHF demonstrator, P2 L/S/C tile, P3 X/Ku/Ka tile and P4 integrated outdoor campaign. Acceptance tests include S-parameters, 3D OTA patterns, radome insertion loss and phase, coherence, synchronization, calibrated AOA, passive-radar Pd/Pfa and resilience after link or power loss.

\begin{figure}[H]
\centering
\includegraphics[width=0.95\textwidth]{fig10_roteiro.png}
\caption{Development and validation roadmap / Roteiro de desenvolvimento e validação.}
\label{fig:roadmap}
\end{figure}

\section*{Português}
A estação requer energia resiliente, UPS, baterias, gerador, enlaces ópticos, gestão térmica, proteção contra descargas, controle de EMC, segurança perimetral e logística. A seleção do sítio deve considerar relevo, horizonte, RFI local, clima, acesso, energia e disponibilidade de fibra. A qualificação ambiental deve abranger vento, vibração, chuva, umidade, salinidade, UV e gradientes térmicos.

A infraestrutura temporal inclui pelo menos duas antenas GNSS externas, com coordenadas e centros de fase levantados. Elas são montadas acima ou próximas ao radome, com visão desobstruída do céu, separação física das Yagis VHF/UHF e cabos filtrados e blindados até a sala temporal. Receptores GNSS multiconstelação redundantes permitem verificação cruzada, detecção de falhas e disciplina do relógio atômico.

A interface civil é uma base de concreto armado em forma de paralelepípedo, com aproximadamente $4\,\mathrm{m}\times4\,\mathrm{m}$ em planta e $3\,\mathrm{m}$ de pé-direito estrutural. Sua laje superior acompanha o plano de corte inferior do radome em $35^\circ$ S, enquanto a área compacta e a elevação maior preservam a visada das antenas nos ângulos inferiores sem apontar sensores para o solo imediato. Há uma abertura de serviço na parede frontal para permitir a entrada de pessoal e equipamentos sem remover os painéis RF.

A solução térmica é setorizada. O calor produzido por LNAs, ADCs, ASICs, FPGAs, transceptores ópticos e conversores DC/DC é conduzido para placas frias locais e rejeitado por caminhos de ar ou líquido monitorados. O núcleo central de sincronismo e fusão possui zona térmica independente. Sensores de condensação, temperatura de entrada e saída e telemetria de potência permitem controle em malha fechada. A deriva térmica de fase é medida, associada ao timestamp e incluída no modelo de calibração.

A solução estrutural utiliza uma moldura interna portante com junções vedadas, placas nodais e reforços. Cada eixo haste/lança da Yagi externa continua do vértice da pirâmide interna através da face triangular, enquanto os painéis dielétricos fornecem proteção ambiental e transmissão RF controlada. A aceitação inclui resistência estática e estabilidade dinâmica: vento, vibração, resposta modal, expansão térmica, chuva, manutenção e ciclos térmicos repetidos.

O caminho de desenvolvimento é P0 simulação de sistema, P1 demonstrador de três nós VHF/UHF, P2 tile L/S/C, P3 tile X/Ku/Ka e P4 campanha externa integrada. Os ensaios incluem S-parâmetros, padrões 3D OTA, perda e fase do radome, coerência, sincronismo, AOA calibrado, Pd/Pfa de radar passivo e resiliência após perda de enlace ou energia.
